\chapter{Corrosion, wear, environmental degradation}
\label{vol10:ch03}

\epigraph{Time and environment are loads; they wear the part the operator never touches.}{}

\chapmeta{Half-life: durable. Archetypes: failure, transport, coupling. Exercise target: 25. Project track: physical.}

\noindent
Corrosion, wear, and environmental degradation are the slow accidents of engineering: components fail not from a single overload but from the cumulative effect of operating environment over time. Corrosion is electrochemical attack; wear is mechanical removal; environmental degradation includes UV damage, thermal cycling, radiation, and chemical attack. The mechanisms are well-understood; the failures persist because the mechanisms are slow, often hidden, and require sustained engineering attention. This chapter develops the engineer's working understanding of corrosion mechanisms at design depth, environmental cracking and hydrogen embrittlement, wear in operating systems, bioenvironmental failure, and the major case studies including the Mianus River Bridge, the Erika oil tanker, and aircraft oxygen-line corrosion.

\section{Corrosion failure mechanisms revisited at depth}\pathtag{standard}

Corrosion is the degradation of a metal through chemical or electrochemical reaction with its environment. The standard mechanisms (covered in Volume V at introductory depth, here at engineering depth):

\textbf{Uniform corrosion:} the metal corrodes uniformly across the surface; the corrosion rate is predictable from material-environment combination tables. Iron in humid air corrodes at $\sim 0.025\,\text{mm/year}$ atmospheric, $\sim 0.1$ to $1\,\text{mm/year}$ in seawater splash zone. Design accommodates the rate with margin (corrosion allowance) or coating; service-life prediction is straightforward.

\textbf{Galvanic corrosion:} two metals in electrical contact in a conductive environment form a battery; the more anodic metal corrodes preferentially. The galvanic series (zinc, aluminum, mild steel, brass, copper, stainless, gold) ranks materials by their potential; the more anodic position (further left) corrodes when coupled to a more cathodic. Practical: aluminum boats with brass fittings corrode at the fittings; carbon-fibre laminates with aluminum fasteners corrode the aluminum. Mitigation: isolating dielectrics, sacrificial anodes (zinc, magnesium), protective coatings.

\textbf{Pitting corrosion:} small areas of localised attack penetrate deeply while the surrounding surface remains intact. Stainless steels in chloride-rich environments are particularly vulnerable; the passive oxide layer breaks down locally; the pit grows with autocatalytic chemistry (the pit interior becomes more aggressive). Pitting can perforate thin sections in years; the bulk of the material remains apparently sound. Detection by visual inspection is unreliable; ultrasonic thickness measurement and pit-depth gauging are standard.

\textbf{Crevice corrosion:} similar to pitting but in occluded geometry (under gaskets, between bolted joints, under deposits). The crevice creates a stagnant electrolyte that becomes aggressive over time. Stainless and aluminum are vulnerable. Design mitigation: avoid crevices, use materials less susceptible (titanium, super-duplex stainless), control the gap (open is better than closed-but-not-tight).

\textbf{Intergranular corrosion:} attack along grain boundaries. Sensitised stainless steels (heated in the $500$-$800\,^\circ\text{C}$ range, depleting chromium at grain boundaries) are vulnerable. The grain boundaries fail; the material remains apparently intact but loses structural integrity. Mitigation: use stabilised stainless (321, 347, with titanium or niobium), avoid welding-heat-affected sensitisation, or post-weld heat-treatment.

\textbf{Selective leaching (dealloying):} one element preferentially dissolves from an alloy. Brass loses zinc (dezincification) in fresh-water service; the brass becomes a porous copper sponge with no strength. Cast iron loses iron (graphitic corrosion) in some soils; the iron leaves the graphite framework behind. Mitigation: inhibited brass alloys, ductile iron replacement.

\textbf{Erosion corrosion:} the combination of mechanical erosion and corrosive attack; flow accelerates the corrosion, particularly when the flow is turbulent (impingement, cavitation). Pump impellers, pipe elbows, valve trim are vulnerable. Mitigation: material upgrade, flow-pattern optimisation, surface protection.

\textbf{Microbiologically influenced corrosion (MIC):} bacteria produce local conditions that drive corrosion. Sulfate-reducing bacteria in anaerobic environments produce $H_2S$ that attacks steel; iron-oxidising bacteria form tubercles that promote pitting beneath. MIC is widespread in marine, oilfield, and water-treatment service. Mitigation: biocides, cleaning, materials less hospitable to bacterial growth.

Corrosion-rate measurement uses weight-loss coupons (long-term), electrochemical methods (linear polarisation resistance, electrochemical impedance spectroscopy, real-time monitoring), and ultrasonic thickness mapping (in-service). The data informs the maintenance plan, the inspection schedule, and the remaining-life estimate.

\begin{archetype}[failure: as time-dependent material loss]
Corrosion and wear are time-dependent material loss. The component is intact at delivery and degraded after years of service. The engineering question is the rate, the location, and the consequence when the loss reaches a critical level. The discipline is recognition (the environment is corrosive or abrasive; this mechanism is a candidate), characterisation (the rate is measured), design (corrosion allowance, coatings, material selection), monitoring (inspection, condition measurement), and renewal (replacement or repair when material loss exceeds limits).
\end{archetype}

\section{Environmental cracking}\pathtag{standard}

Environmental cracking is the brittle failure of a normally-ductile material due to the combined action of stress and environment. The standard mechanisms:

\textbf{Stress corrosion cracking (SCC):} the metal fails by brittle cracking under tensile stress in a specific environment, at stresses well below the yield strength. The crack propagates along grain boundaries (intergranular SCC) or through grains (transgranular SCC). Examples: austenitic stainless steel in chloride environments (Cl-SCC); brass in ammonia environments; aluminum alloys in chloride environments; high-strength steels in hydrogen-bearing environments.

The SCC mechanism involves the simultaneous action of:
1. \textbf{Susceptible material:} not all metal-environment combinations crack; the specific pairing matters.
2. \textbf{Tensile stress:} compressive stress does not produce SCC. Residual stresses from welding, cold-working, or assembly are sufficient.
3. \textbf{Specific environment:} the cracking chemistry is specific; small changes in concentration, temperature, or pH can switch cracking on or off.

The SCC time-to-cracking depends on stress and environment; under sustained stress in the cracking environment, failures occur in weeks to years. The unpredictability and the brittle nature make SCC particularly dangerous.

Mitigation: material selection (e.g., duplex stainless in chloride environments rather than 304/316), stress-relief annealing of welds, surface treatment (shot peening introduces compressive residual stress), inhibitors in the environment, monitoring and inspection.

\textbf{Liquid-metal embrittlement (LME):} a normally-ductile metal becomes brittle when in contact with a specific liquid metal. Aluminum embrittled by mercury, steel by liquid zinc, copper by lead. Mitigation: avoid contact between the specific metal pairs.

\textbf{Solid-metal embrittlement:} similar to LME but at lower temperatures with metals that diffuse along grain boundaries. Less common but documented for specific pairings.

\textbf{Corrosion fatigue:} cyclic stress in a corrosive environment produces faster crack growth than either alone. The fatigue limit disappears (every cycle contributes some damage); the safe-life prediction shortens significantly. Mitigation: avoid the cracking environment, surface treatment, fatigue-resistant material selection.

\section{Hydrogen embrittlement}\pathtag{standard}

Hydrogen embrittlement is a specific environmental-cracking mechanism in which atomic hydrogen diffuses into the metal and reduces its ductility, leading to brittle failure. The mechanism affects high-strength steels (especially above $1000\,\text{MPa}$ yield strength), titanium alloys, and some nickel alloys.

Sources of hydrogen:
\textbf{Cathodic charging:} the metal is the cathode in an electrochemical process (cathodic protection, electroplating, pickling); hydrogen is generated at the surface and absorbed.
\textbf{Service environment:} the operating environment contains $H_2$, $H_2S$, or other hydrogen sources. Oil and gas service ("sour service" with $H_2S$) is the major industrial example.
\textbf{Welding:} the weld process introduces hydrogen from moisture in the electrode coating or in the gas; the heat-affected zone is vulnerable.
\textbf{Manufacturing:} acid pickling, electroplating, and similar processes introduce hydrogen.

The hydrogen-embrittlement failures are characteristically brittle (no visible plastic deformation), delayed (failures may occur hours to days after the hydrogen exposure), and at stresses below the yield strength. The fracture surface shows brittle features (cleavage, intergranular).

Standard tests: NACE TM0177 (slow-strain-rate test in $H_2S$ environment), ASTM F1624 (hydrogen-embrittlement-prone-to-fracture threshold). The tests qualify a material-environment combination for the intended service.

Mitigation: baking (heating after hydrogen-charging operations to drive out hydrogen), avoiding high-strength steels in hydrogen environments, using hydrogen-resistant materials (low-strength steels, austenitic stainless with chromium content adjusted), surface coatings, electrochemical treatment.

\section{Wear in operating systems}\pathtag{standard}

Wear is the mechanical removal of material from surfaces in relative motion. The dominant mechanisms (covered briefly in Vol VIII Ch 2, here as failure mode):

\textbf{Abrasive wear:} hard particles or asperities cut the softer surface. Sand in fluid systems, dust in air systems, hard particles in lubricant. Mitigation: filtration, hard surface treatment, abrasive-resistant materials.

\textbf{Adhesive wear:} the surfaces stick locally and tear; metal transfers between them. Galling (severe adhesive wear) in poorly-lubricated metal-on-metal contact, particularly stainless-on-stainless. Mitigation: dissimilar materials, lubrication, surface treatment.

\textbf{Erosive wear:} fluid-borne particles impinge on the surface and remove material. Slurry pumps, sand-laden flow, particulate exhaust. Mitigation: surface hardness, flow-pattern optimisation.

\textbf{Cavitation wear:} bubbles in a liquid collapse near a surface, producing local shock that removes material. Pump impellers, ship propellers, valve trim in throttling service. Mitigation: cavitation-resistant materials (some stainless grades, Stellite), avoidance of cavitation conditions.

\textbf{Fretting wear:} small-amplitude oscillatory motion at apparently stationary contacts (bolted joints, press fits, splined connections). The motion is below visual detection but produces local wear and fatigue. Mitigation: lubrication, surface treatment, stiff joint design.

\textbf{Fatigue wear (surface fatigue):} repeated contact stress produces subsurface fatigue cracks that grow and produce spalling. Rolling-element bearings, gears at the tooth surface. Mitigation: material cleanliness, surface treatment, load reduction.

Wear rates are commonly $10^{-12}$ to $10^{-15}\,\text{m}^3/(\text{N}\cdot\text{m})$ for lubricated contacts; orders of magnitude higher for unlubricated. The Archard wear equation: $V = K \cdot F \cdot L / H$, with $V$ the wear volume, $K$ a wear coefficient, $F$ the load, $L$ the sliding distance, $H$ the hardness. The equation is approximate but useful for ranking material combinations.

Wear monitoring uses oil analysis (wear-debris characterisation), vibration analysis (changes in bearing or gear signatures), and direct measurement (dimensional checks, profilometry). Predictive maintenance based on wear-debris monitoring is mature for industrial equipment.

\section{Bioenvironmental failure (medical implants, marine systems)}\pathtag{mastery}

Bioenvironmental failure addresses materials in biological environments: medical implants in the human body, marine structures in seawater with biological growth, food-processing equipment, agricultural systems.

Medical implants face specific challenges:

\textbf{Biocompatibility:} the implant must not provoke adverse reactions in the host. The standards (ISO 10993) define biological testing. The body's response includes inflammation, fibrous encapsulation, possibly rejection.

\textbf{Corrosion in body fluids:} the body environment is corrosive (saline, $37\,^\circ\text{C}$, biological molecules). Stainless steel (316LVM), cobalt-chromium, and titanium alloys are standard. Failure modes include pitting, crevice corrosion, and (for high-strength steels) hydrogen embrittlement from biological cathodic processes.

\textbf{Wear:} hip and knee replacements are wear-loaded; the wear debris is a major failure mode (induces osteolysis, loosening). Modern designs use highly-cross-linked polyethylene against ceramic or metal counter-faces; wear rates are an order of magnitude below older designs.

\textbf{Fatigue:} cyclic load (millions of cycles over an implant lifetime) drives fatigue failure of cardiovascular stents, dental implants, orthopaedic devices. Material toughness, surface condition, and design detail are critical.

\textbf{Combined modes:} corrosion fatigue, fretting corrosion, and stress-corrosion cracking are documented failure modes in medical implants.

Marine systems face their own challenges:

\textbf{Biofouling:} barnacles, mussels, algae, and bacteria colonise immersed surfaces; the biofilm and macrofouling increase drag, weight, and corrosion. Antifouling coatings (copper-based, silicone-based, biocide-eluting) mitigate; the field continues to evolve.

\textbf{Microbiologically influenced corrosion:} as discussed in the corrosion section. Marine environments are particularly rich in microorganisms.

\textbf{Coupled corrosion-fatigue:} marine structures (offshore platforms, ship hulls) experience cyclic wave loading in corrosive seawater. The fatigue life is reduced; the inspection programme must address the combination.

\textbf{Cathodic protection:} sacrificial anodes (zinc, aluminum) or impressed-current systems protect against corrosion. The design of the cathodic protection is engineering work; under-protection allows corrosion; over-protection produces hydrogen embrittlement of high-strength steels.

\section{Case: Mianus River Bridge, Erika, Boeing 707 oxygen-line corrosion}\pathtag{core}

Three corrosion case studies illustrate the discipline's range:

\textbf{Mianus River Bridge (1983):} a section of I-95 in Connecticut collapsed, killing three. The investigation found corrosion of the pin-and-hanger suspension system: water had penetrated the pin clevis, corroded the steel, and weakened the connection until it failed under traffic loads. The corrosion was hidden: the connection was not easily inspectable; the bridge inspection had not penetrated the protected geometry. The failure prompted the development of pin-and-hanger replacement programs across US bridges and accelerated the move to redundant load-path designs.

\textbf{Erika oil tanker (1999):} a $25$-year-old single-hull tanker broke in half off the French coast, spilling $19{,}800$ tonnes of fuel oil. The investigation identified extensive corrosion in the ballast tanks (not the cargo tanks); the structural steel had thinned to a fraction of its original thickness; the cyclic loading from waves and cargo handling overcame the weakened structure. The failure accelerated the EU phase-out of single-hull tankers and tightened inspection requirements for aged tankers.

\textbf{Boeing 707 oxygen-line corrosion (1990s):} a series of in-flight oxygen-line ruptures on aged 707 aircraft was traced to corrosion of the stainless-steel oxygen distribution lines. The corrosion mechanism: oxygen-line cleanliness procedures used chemicals that left chloride residues; over decades, the residues attacked the stainless steel via Cl-SCC. The failures were intermittent but potentially catastrophic (oxygen fires). The fix involved cleaning procedure changes, line replacement programs, and enhanced inspection. The case illustrates how a procedural-cleanliness shortcoming can produce structural failure modes years later.

The common pattern: corrosion is slow, often hidden, and requires sustained engineering attention. The failures are not at the design point but at the integration of design, operation, environment, and time.

\begin{principle}[The corrosion-and-wear-allowance principle]
For components in corrosive or wear-prone service, the engineering design must include explicit allowance for material loss over the design life. The allowance includes: corrosion-rate estimation (from material-environment data), corrosion-allowance thickness (material to lose during the design life), inspection access (geometry that allows condition measurement), monitoring (sensors or periodic inspection), and replacement criteria (when remaining material approaches the structural minimum). Designs without allowance produce premature failures; designs with allowance and disciplined maintenance can serve for decades.
\end{principle}

\section{Project}\pathtag{core}

\begin{project}[Assess corrosion risk for one infrastructure asset]
\textbf{Objective.} Assess the corrosion risk for one infrastructure asset in the reader's environment. Identify the dominant corrosion mechanisms; estimate the rate; recommend monitoring approach.

\textbf{Method.} Choose an asset: a steel water pipe (residential or municipal), a bridge in your area, an outdoor steel structure (utility tower, fence), a building's HVAC system, an automotive component. Identify the metal, the environment (atmosphere, water, soil), the geometric features that affect corrosion. Estimate the dominant corrosion mechanism (uniform, galvanic, pitting, crevice, MIC). Estimate the corrosion rate from published material-environment data or your own measurements. Recommend a monitoring approach: visual inspection frequency, ultrasonic thickness measurement, coupon installation, or other appropriate method.

\textbf{Deliverables.} The asset description and environment characterisation; the corrosion-mechanism identification with rationale; the rate estimate with sources; the monitoring recommendation; a written analysis of the trade-off between monitoring cost and the consequence of undetected failure.

\textbf{Hazard.} Field inspection should be conducted safely with appropriate permission and PPE. Some infrastructure inspection requires professional licensure.
\end{project}

\section{Exercises}

\subsection*{Calculation}\pathtag{core}

\begin{exercise}
A steel pipe is exposed to atmospheric corrosion at $0.05\,\text{mm/year}$. The original wall thickness is $6\,\text{mm}$; the structural minimum is $3\,\text{mm}$. Compute the remaining service life.
\end{exercise}

\begin{exercise}
Two coupled metals form a galvanic couple in seawater: brass and steel. State which corrodes and qualitatively the rate change relative to the steel alone.
\end{exercise}

\begin{exercise}
A steel structure has corroded to $80\%$ of original thickness after $20$ years. The corrosion allowance was $1\,\text{mm}$ over the original $5\,\text{mm}$ thickness. State whether the structure has exceeded its corrosion budget.
\end{exercise}

\begin{exercise}
A bearing has wear rate $10^{-14}\,\text{m}^3/(\text{N}\cdot\text{m})$; load $1000\,\text{N}$, sliding distance $10^7\,\text{m/year}$. Compute the annual wear volume.
\end{exercise}

\begin{exercise}
A pump impeller in a cavitating service shows $0.5\,\text{mm}$ material loss per $1000$ hours. State the expected service life if structural minimum allows $2\,\text{mm}$ loss.
\end{exercise}

\begin{exercise}
A medical implant of titanium alloy is exposed to body fluids; reported corrosion rate is $0.001\,\text{mm/year}$. Compute the expected mass loss over $30$ years if the surface area is $50\,\text{cm}^2$ and density $4.5\,\text{g/cm}^3$.
\end{exercise}

\subsection*{Derivation}\pathtag{standard}

\begin{exercise}
Derive the Archard wear equation from first principles: real-area-of-contact and material-removal models.
\end{exercise}

\begin{exercise}
Derive the relationship between cathodic current density and protection in cathodic-protection systems.
\end{exercise}

\begin{exercise}
Derive the Faraday-based corrosion rate from electrochemical current density.
\end{exercise}

\subsection*{Estimation}\pathtag{standard}

\begin{exercise}
Estimate the corrosion rate of carbon steel exposed to: (a) dry indoor atmosphere; (b) outdoor coastal atmosphere; (c) seawater splash zone; (d) buried in clay soil.
\end{exercise}

\begin{exercise}
Estimate the inspection-cost-vs-replacement-cost trade-off for a buried steel pipe: continuous monitoring vs scheduled excavation inspection vs replace-at-end-of-life.
\end{exercise}

\begin{exercise}
Estimate the time to SCC failure for a stressed austenitic stainless component in a moderately chloride-rich environment.
\end{exercise}

\subsection*{Simulation}\pathtag{standard}

\begin{exercise}
Simulate the remaining life distribution for a corroding pipe given uncertainty in corrosion rate and inspection accuracy.
\end{exercise}

\begin{exercise}
Simulate the wear-particle generation in a lubricated bearing under increasing load; identify the transition to failure.
\end{exercise}

\subsection*{Design}\pathtag{standard}

\begin{exercise}
Design a corrosion-protection strategy for a buried steel pipeline: coating, cathodic protection, monitoring, replacement plan.
\end{exercise}

\begin{exercise}
Design a wear-monitoring program for a fleet of industrial pumps: oil analysis, vibration monitoring, performance trending.
\end{exercise}

\begin{exercise}
Design the materials specification for an offshore platform's seawater systems to avoid Cl-SCC and crevice corrosion.
\end{exercise}

\subsection*{Judgement}\pathtag{core}

\begin{exercise}
A team's offshore platform shows accelerated corrosion in ballast tanks not predicted by the design. State three engineering hypotheses.
\end{exercise}

\begin{exercise}
A team is choosing between a corrosion-resistant material (more expensive) and a corrosion-allowance approach (cheaper material with thicker section) for a chemical-plant vessel. State the engineering trade-offs.
\end{exercise}

\begin{exercise}
A team's medical implant has failed in service after $5$ years; the design life was $20$. State the engineering investigation including possible mechanisms.
\end{exercise}

\begin{exercise}
A team's pump impeller wear is significantly faster than predicted. State the engineering investigation.
\end{exercise}

\begin{exercise}
A team is asked to certify a bridge for $50$ more years of service; the original design life was $50$ years and the bridge is $60$ years old. State the engineering considerations.
\end{exercise}

\begin{exercise}
A team's stainless-steel storage tank shows pitting after $2$ years in a service the manufacturer claimed compatible. State the engineering response.
\end{exercise}

\begin{exercise}
A team is investigating an unexpected SCC failure in an aluminum aerospace component. State the engineering investigation including environmental and stress history.
\end{exercise}

\begin{exercise}
A team's marine structure has both fouling and corrosion concerns. State the engineering response and the interaction between the two.
\end{exercise}
